%%%%%%%%%%  scan idx 454  |  folio 441  %%%%%%%%%%
\textbf{Proposition 3.4.} Avec les hypothèses et les notations de 3.1 et 3.3, on a pour tout \(x\in K^\circ(Y)\) la relation:
\[
f^*(i_*(x))=j_*\bigl(g^*(x).\lambda_{-1}(F)\bigr).
\]

Comme les deux membres de la relation à démontrer sont additifs, on peut supposer que \(x\) est la classe d'un \(\cO_Y\)-Module localement libre de type fini \(M\). L'élément \(f^*(i_*(x))\) est alors la classe du complexe \(\mathbb Lf^*(\mathbb Ri_*(M))\), obtenu en prenant une résolution projective de \(i_*(M)\) et en lui appliquant le foncteur \(f_*\). Les faisceaux d'homologie du complexe ainsi obtenu sont les \(\underline{\Tor}_i^{\cO_X}(M,\cO_{X'})\). Or d'après 2.4 et 3.2, on a:
\[
\underline{\Tor}_i^{\cO_X}(M,\cO_{X'})
\xrightarrow{\sim}
M\otimes_{\cO_Y}\underline{\Tor}_i^{\cO_X}(\cO_{X'},\cO_Y)
\xrightarrow{\sim}
M\otimes_{\cO_Y}\bigwedge^iF.
\]

Il en résulte que ces faisceaux d'homologie sont localement libres de type fini sur \(\cO_{Y'}\), donc parfaits sur \(\cO_{X'}\); par suite, la classe dans \(K^\circ(X')\) du complexe \(\mathbb Lf^*(\mathbb Ri_*(M))\) est la somme alternée des classes de ses faisceaux d'homologie, soit:
\[
f^*(i_*(x))
=\sum_{i\ge0}(-1)^i\operatorname{cl}
\bigl(j_*(M\otimes_{\cO_Y}\bigwedge^iF)\bigr)
=\sum_{i\ge0}(-1)^i
j_*\bigl(g^*(x)\lambda^i(F)\bigr).
\]

\textbf{Lemme 3.5.} Avec les hypothèses et les notations de 3.1, on a les isomorphismes suivants :
\begin{enumerate}[label=\roman*)]
\item pour tout \(n\ge0\), et pour tout \(i\ge1\), \(R^if_*\bigl(\cO_{X'}(n)\bigr)=0\);
\item pour tout \(n\ge0\), \(f_*\bigl(\cO_{X'}(n)\bigr)=J^n\).
\end{enumerate}

On peut évidemment supposer \(X\) affine, donc quasi-compact.

i) Rappelons d'abord le lemme suivant, qui a été démontré dans exp. III.

